\section{方式一：AR直接跟随（第二小组）}
\subsection{实验目的与方案}
本组使用Kinect识别前车背部的ALVAR标记，根据标签与后车的相对位置调整速度和方向。我们希望后车既能保持一定车距，又能在前车转弯时继续对准目标，并在扫描或目标信息异常时停止。\cite{follow1}

实验中使用了两种直接跟随程序：\code{follow_ar.py}直接在相机坐标下计算速度，便于检查识别与底盘响应；\code{marker_follower.py}先将标记转换到车体坐标，再通过独立门控输出速度。两种程序的距离参数和保护条件不同，调试时分别设置，表中列出主要区别。

\begin{table}[htbp]\centering\small
\caption{两种AR直接跟随程序的比较}
\begin{tabularx}{\linewidth}{p{31mm}YY}\toprule
项目 & 相机坐标直接控制 & 车体坐标与独立门控\\\midrule
功能包 & \code{ar_follower} & \code{ar_marker_following}\\
目标参数 & \code{target_id}，默认1 & \code{marker_id}，默认1\\
跟随距离 & 默认0.4\,m & 默认1.0\,m\\
位置坐标 & 光学坐标$z$前向、$x$横向 & TF到车体坐标$x$前向、$y$左向\\
速度限制 & 线速度0.3\,m/s、角速度1.0\,rad/s & 前进上限0.12\,m/s、角速度0.45\,rad/s\\
输出 & 直接发\code{teleop}输入 & 通过门控后发导航输入\\
停止条件 & 空检测停车，时效保护不完整 & 检查目标、任务、扫描及指令时效\\
障碍响应 & 停止前进并慢速转向 & 门控停车，不负责路径绕障\\\bottomrule
\end{tabularx}\end{table}

课程还介绍了固定场地标签航路点的方法。我们这里选择把标签贴在前车上，利用相对位姿跟随，因此不需要后车预先知道前车的全局位置；相应地，标签的可见性对运行效果影响较大。

\subsection{标记识别、尺寸与目标ID}
\code{ar_track_alvar}接收RGB图像与CameraInfo，输出\code{AlvarMarkers}。跟随节点只选取参数指定ID，识别节点能看见其他ID也不代表控制节点会跟随。改变目标时，旧节点使用\code{_target_id:=1}，新版launch使用\code{marker_id:=1}，应核对实际入口后设置。

标签应为ALVAR格式，实际黑边尺寸须与\code{marker_size_cm}一致。旧实验launch默认14\,cm，新包识别launch默认9\,cm，均不是现场尺寸已确定的结论。尺寸不一致会引入尺度误差；标记反光、视野边缘、遮挡、打印缩放和相机内参错误也会影响距离估计。多车应为直接前车分配明确标记ID，并分别登记图像、内参话题与光学坐标系。

\subsection{相机直接控制与车体坐标控制}
早期节点直接使用相机光学坐标$z_m,x_m$，控制思想为
\begin{equation}
v=\operatorname{sat}\!\left(K_v(z_m-d_0)\right),\qquad
\omega=\operatorname{sat}\!\left(-K_\omega x_m\right).
\end{equation}
光学坐标中横向向右为正，ROS角速度向左为正，故转向项取负号。实机需核对坐标约定与安装方向，不能在TF或相机配置错误时仅用反号掩盖问题。旧节点扫描前方60度扇区，小于0.25\,m时线速度归零、角速度设为0.5\,rad/s；目标太近时也停止前进，但扫描失联没有独立超时保护。

新\code{marker_follower.py}将量测在原始时间戳变换到\code{base_frame}，保留有限且位于前方的目标$(x,y)$。用$\alpha=\operatorname{atan2}(y,x)$计算方位，源码控制律为
\begin{align}
v&=\begin{cases}\min\{0.12,\,0.4(x-d_0)\},&x>d_0\land|\alpha|<0.4,\\0,&\text{其他情况},\end{cases}\\
\omega&=\begin{cases}\operatorname{clamp}(1.2\alpha,0.45),&x>0.65,\\0,&\text{其他情况}.\end{cases}
\end{align}
距离以米、角度以弧度计。新版观测输入须在0.5\,s内，控制阶段允许的最近目标年龄为0.6\,s，任务状态年龄为1.5\,s；只有carrying、releasing、done且未急停时发布有效许可。新版本不能依赖旧版“直接取相机$z$”来解释车体$x$的数值。

\subsection{独立速度门控与停止条件}
新版原始速度发布到\code{cmd_vel_raw}，许可发布到\code{motion_allowed}，再由\code{velocity_gate.py}检查并输出\code{cmd_vel_safe}。任务状态由上游控制器转换为许可，门控本身不直接订阅任务字符串。许可或速度超过0.4\,s未更新、扫描超过0.5\,s未刷新、急停有效、速度非有限或为倒车时输出零速度。

默认门控最大线速度0.15\,m/s、最大角速度0.5\,rad/s，前方停止距离0.5\,m；直接跟随控制器自身限制更低，不能用门控上限替代实际输出上限。直行扫描半扇区为45度，转动时扩大到90度，至少需要足够有效量测；NaN/Inf缺测不能当作安全。壁钟线程用于处理ROS时钟暂停。目标丢失仍可能在目标时效窗口内沿用最近有效观测，不能写成瞬时硬件急停。

\subsection{调试流程与启动检查}
先确认底盘、相机、实际图像话题和标签ID，再检查TF与\code{/scan}；只运行识别节点时不应期待底盘运动。任务状态、许可或扫描任一条件未满足，门控保持停止是预期行为。多车必须统一\code{robot_ns}，识别launch默认\code{tb2}、跟随launch默认\code{tb1}，不能同时依赖两套默认命名空间。

程序分成多个功能包后，启动文件中仍有部分节点保留了原包名。运行前需要把launch中的包名与脚本实际位置对应起来，再重新编译和检查节点。这个问题提醒我们，修改程序组织时不能只移动文件，还要一起检查启动配置和依赖。

\subsection{队列扩展与方式独立性}
可令TonyPi贴ID0、首车跟随ID0，首车贴ID1、次车跟随ID1。每台车拥有唯一命名空间和TF，但改变ID只改变目标，不改变直接跟随的控制律。该方案不提供地图路径规划，遇障会停车而非主动绕行；需要绕障时，应接入方式二导航或由对应组明确其他规划实现。

\subsection{实验现象与改进}
单标记调试主要检查标签移动与坐标变化的对应关系，以及距离和方位误差怎样影响速度。标签偏向一侧时，应产生相应转向；距离接近期望值时，应逐渐减速。目标ID切换和多车场景还需要分别运行，不能只改参数后就认为整个队列已经调好。

离线检查通过了11项控制与停止逻辑测试。它们有助于检查条件分支，但不能代替实际相机、TF和底盘运行。后续还需要记录不同距离下的误差、转弯时的丢失时长，以及扫描或任务状态中断后的停车情况。

\subsection{小组合作与体会}
陈凯华负责标签、相机和坐标检查，刘彦廷负责多机配置与场景，\LuoMin{}负责控制和门控，许程奕负责试验记录与分析。我们共同进行参数调整和现场测试，具体分工见第\ref{sec:division}节。调试中最有帮助的方法是同时查看识别坐标和速度输出，这样能够区分“没有检测到目标”和“控制条件没有满足”。
